\begin{figure}[ht]
\centering
\begin{tikzpicture}[x=1cm, y=1cm, font=\small]

% Title
\node[anchor=south, font=\bfseries] at (5.5, 5.7) {Volume by water displacement (Archimedes)};

% Panel A: graduated cylinder before immersion
\begin{scope}[shift={(0,0)}]
  \node[anchor=south, font=\small\itshape] at (1.5, 5.0) {(a) before};
  % Cylinder body
  \draw[thick] (0.5, 0.4) -- (0.5, 4.8) (2.5, 0.4) -- (2.5, 4.8);
  \draw[thick] (0.5, 0.4) arc[start angle=180, end angle=360, x radius=1, y radius=0.25];
  \draw[thick, densely dashed] (0.5, 0.4) arc[start angle=180, end angle=0, x radius=1, y radius=0.25];
  % Water fill to level V0
  \fill[engblue!25] (0.5, 0.4) rectangle (2.5, 2.4);
  \draw[thick, engblue] (0.5, 2.4) -- (2.5, 2.4);
  \draw[thick, engblue] (0.5, 2.4) arc[start angle=180, end angle=0, x radius=1, y radius=0.18];
  % Graduations
  \foreach \y in {1.0, 1.5, 2.0, 2.5, 3.0, 3.5, 4.0} {
    \draw[thick] (2.4, \y) -- (2.5, \y);
  }
  % Level label
  \draw[->, thick, engred] (3.4, 2.4) -- (2.6, 2.4);
  \node[anchor=west, font=\small] at (3.4, 2.4) {$V_{0}$};
\end{scope}

% Panel B: graduated cylinder after immersion with stone
\begin{scope}[shift={(6,0)}]
  \node[anchor=south, font=\small\itshape] at (1.5, 5.0) {(b) after};
  % Cylinder
  \draw[thick] (0.5, 0.4) -- (0.5, 4.8) (2.5, 0.4) -- (2.5, 4.8);
  \draw[thick] (0.5, 0.4) arc[start angle=180, end angle=360, x radius=1, y radius=0.25];
  \draw[thick, densely dashed] (0.5, 0.4) arc[start angle=180, end angle=0, x radius=1, y radius=0.25];
  % Water to V1 (higher than V0)
  \fill[engblue!25] (0.5, 0.4) rectangle (2.5, 3.4);
  \draw[thick, engblue] (0.5, 3.4) -- (2.5, 3.4);
  \draw[thick, engblue] (0.5, 3.4) arc[start angle=180, end angle=0, x radius=1, y radius=0.18];
  % Stone immersed
  \fill[engblack!70] plot[smooth cycle, tension=0.7] coordinates {(1.0, 1.4) (1.4, 1.1) (2.0, 1.3) (2.1, 1.8) (1.8, 2.2) (1.3, 2.1) (0.9, 1.8)};
  % Graduations
  \foreach \y in {1.0, 1.5, 2.0, 2.5, 3.0, 3.5, 4.0} {
    \draw[thick] (2.4, \y) -- (2.5, \y);
  }
  % Level label
  \draw[->, thick, engred] (3.4, 3.4) -- (2.6, 3.4);
  \node[anchor=west, font=\small] at (3.4, 3.4) {$V_{1}$};
  % Stone label
  \node[anchor=west, font=\scriptsize] at (2.3, 1.6) {object};
\end{scope}

% Bottom equation
\node[anchor=north, font=\small] at (5.5, -0.1) {$V_{\text{object}} = V_{1} - V_{0}$};
\node[anchor=north, font=\scriptsize] at (5.5, -0.7) {valid only if the object is fully wetted, non-absorbent, and free of trapped air bubbles};

\end{tikzpicture}
\caption[Volume by water displacement]{Volume by water displacement.
The object's volume equals the rise in water level between the
before-immersion and after-immersion readings. The method assumes the
object is fully wetted, free of trapped air bubbles, and neither
soluble nor absorbent in water. The reading's uncertainty has three
sources: the cylinder's graduation interval (resolution), the
parallax of the meniscus reading, and the volume of water that
adheres to the object on removal. Section 7.5 returns to displacement
as one method in the chapter's three-method volume project. Source:
editors' schematic.}
\label{fig:vol01:ch07:volume-displacement}
\end{figure}
